\documentclass[10pt,a4paper]{article}

% Two-page submission version; compile with LuaLaTeX.
\usepackage[margin=0.68in,headheight=13pt,headsep=10pt,footskip=18pt]{geometry}
\usepackage[default]{sourcesanspro}
\usepackage{amsmath,booktabs,microtype,titlesec,fancyhdr}
\usepackage[font=footnotesize,labelfont=bf,skip=3pt]{caption}
\usepackage[hidelinks]{hyperref}

\setlength{\parindent}{0pt}
\setlength{\parskip}{3pt}
\raggedbottom
\titleformat{\section}{\fontsize{12}{14}\selectfont\bfseries}{\thesection}{0.55em}{}
\titlespacing*{\section}{0pt}{5pt}{2pt}
\setlength{\textfloatsep}{5pt plus 2pt minus 2pt}
\setlength{\floatsep}{4pt plus 2pt minus 2pt}

\pagestyle{fancy}
\fancyhf{}
\renewcommand{\headrulewidth}{0pt}
\fancyhead[L]{\fontsize{8}{9}\selectfont EPFL CS-433 Machine Learning}
\fancyhead[R]{\fontsize{8}{9}\selectfont Project 1 -- Team T123}
\fancyfoot[C]{\fontsize{8}{9}\selectfont \thepage}

\begin{document}

{\centering
{\fontsize{15}{17}\selectfont\bfseries Predicting Cardiovascular Disease Risk from BRFSS Data\par}
\vspace{3pt}
{\small Haobin Wang, Yunxin Huang, Chaoqi Xu \hfill Team T123, Autumn 2026\par}
}

\section{Introduction and Dataset}

We study binary classification of myocardial infarction or coronary heart
disease (\texttt{\_MICHD}) using the 2015 Behavioral Risk Factor Surveillance
System (BRFSS). The dataset contains 328,135 labeled respondents and 321 health
and lifestyle features: 28,975 positive cases and 299,160 negative cases, so
only 8.83\% of observations are positive. Consequently, ordinary accuracy can
favor the negative class. Our objective is to compare
L2-regularized logistic regression, a class-weighted linear SVM, and a random
forest, and to assess how preprocessing, imbalance handling, feature
engineering, and decision-threshold selection affect positive-class detection.
The target was encoded as $+1$ for disease and $-1$ otherwise. In accordance
with the project rules, model training and evaluation were implemented with
NumPy rather than external machine-learning libraries.

\section{Preprocessing and Validation}

We used the same 80/20 stratified training--validation split with
\textbf{seed 42} for all models, giving 262,508 training and 65,627 validation
examples while preserving class proportions. Every
preprocessing parameter was learned on the training partition and applied
unchanged to validation data. The \emph{Basic} pipeline removed constant and
entirely missing columns, median-imputed missing values, and standardized the
remaining variables. The \emph{BRFSS} pipeline additionally used the survey
codebook to map feature-specific ``Don't know'' and ``Refused'' codes to
missing, added numerical missingness indicators, and one-hot encoded
categorical variables. We report accuracy, balanced accuracy, precision,
recall, and positive-class F1. F1 summarizes the precision--recall compromise,
whereas balanced accuracy gives equal importance to sensitivity and
specificity; both are more informative than raw accuracy under this imbalance.
Specifically, $F1=2PR/(P+R)$ and balanced accuracy is the mean of recall and
specificity. Test IDs were kept separate from features so predictions could be
returned in the original AIcrowd order.

\section{Models}

\textbf{3.1 L2-regularized logistic regression.}
Basic preprocessing retained 315 of 321 features. All 23,180 positive training
examples were paired with an equal random sample of negatives. Our NumPy model
minimized binary cross-entropy plus $\lambda\lVert\mathbf{w}\rVert_2^2$
without penalizing the intercept. Mini-batch gradient descent used 20 epochs,
batch size 2,048, $\lambda=10^{-4}$, and learning rate
$0.1/\sqrt{1+e/5}$. Thresholds maximizing accuracy, balanced accuracy, and F1
were compared on the unchanged validation set. The fitted probability score
$\sigma(\mathbf{x}^{\top}\mathbf{w}+b)$ was converted to a class label only
after selecting the operating threshold, leaving the learned weights fixed.
Accuracy and balanced-accuracy thresholds were searched on a 0.005 grid; the
F1 threshold was selected from the distinct validation scores.

\textbf{3.2 Class-weighted linear SVM.}
The SVM used BRFSS preprocessing, yielding 1,194 features including the
intercept. Its NumPy soft-margin objective was
$N^{-1}\sum_i\alpha_i\max(0,1-y_i\mathbf{w}^{\top}\mathbf{x}_i)
+\frac{\lambda}{2}\lVert\mathbf{w}\rVert_2^2$, with positive weight 10.3248
and negative weight 1. Stochastic subgradient descent used
$\lambda=0.01$, learning rate $10^{-4}$, 200,000 updates, and seed 42.
Validation F1 selected decision threshold 0.5009 instead of the standard
threshold zero. We evaluated 302 candidate thresholds derived from validation
decision scores, using balanced accuracy only to break an F1 tie.

\textbf{3.3 Random forest.}
Basic preprocessing outperformed the high-dimensional BRFSS representation.
The forest averaged positive-class probabilities from 60 bootstrap-trained
trees using Gini impurity, maximum depth 12, and
$\lfloor\sqrt{d}\rfloor$ candidate features per split. A functional-limitation
count derived from five survey items raised validation F1 from 0.417 to 0.424.
Invalid responses in these items were treated as missing. Thresholds from 0.15
to 0.25 were evaluated in increments of 0.01, selecting 0.18; greater depth,
more candidate features, and class resampling gave negligible gains at higher
computational cost. Unlike the two linear models, the forest can represent
nonlinear interactions without explicit polynomial feature expansion.

\newpage

\section{Results}

\begin{table}[h]
\centering
\small
\caption{Local validation performance at each model's F1-selected threshold.}
\label{tab:comparison}
\begin{tabular}{lrrrrrr}
\toprule
Model & Threshold & Accuracy & Balanced acc. & Precision & Recall & F1 \\
\midrule
L2 logistic regression & 0.737 & \textbf{0.871} & 0.711 & 0.345 & 0.516 & 0.414 \\
Weighted linear SVM & 0.5009 & 0.868 & 0.732 & \textbf{0.349} & 0.566 & \textbf{0.432} \\
Random forest & 0.18 & 0.856 & \textbf{0.737} & 0.325 & \textbf{0.593} & 0.424 \\
\bottomrule
\end{tabular}
\end{table}

Threshold selection substantially changed minority-class detection. Logistic
regression's accuracy-oriented threshold 0.940 produced accuracy 0.913 but only
F1 0.110 and recall 0.061; its F1 threshold 0.737 increased F1 to 0.414 and
recall to 0.516. For the random forest, threshold 0.5 similarly produced
accuracy 0.912 but only F1 0.015 and recall 0.008, whereas threshold 0.18 raised
F1 to approximately 0.424 and recall to 0.593.

For the SVM, moving the threshold from 0 to 0.5009 reduced false positives from
11,147 to 6,126 and raised precision from 0.275 to 0.349. Local accuracy rose
from 0.806 to 0.868 and F1 from 0.399 to 0.432, while recall decreased from
0.729 to 0.566 and balanced accuracy from 0.771 to 0.732. Thus, threshold 0.5009
is F1-oriented rather than uniformly superior. On AIcrowd, it obtained
\textbf{F1 0.431 and accuracy 0.866}, closely matching local validation
(0.432 and 0.868). The earlier threshold-zero submission obtained F1 0.401 and
accuracy 0.804.

The all-negative reference illustrates why raw accuracy is insufficient: it
achieves 91.17\% accuracy but zero recall and F1 for the positive class. Across
the tuned models, the SVM achieved the highest local F1, the forest achieved
the highest recall and balanced accuracy, and logistic regression achieved the
highest ordinary accuracy.

\section{Discussion}

The results show that the operating threshold is as important as the model
family under strong imbalance. Default or accuracy-oriented thresholds favored
the majority class, whereas F1-based calibration produced useful positive-class
detection. The SVM offered the strongest F1 compromise: raising its threshold
removed many false positives and improved precision and accuracy, but inevitably
missed more positive cases. The preferred threshold therefore depends on the
relative cost of false alarms and missed cardiovascular-risk cases. The close
agreement between the SVM's local and AIcrowd scores suggests that the shared
holdout was representative of the hidden test set.

The comparison does not isolate model family completely. Logistic regression
used balanced sampling with Basic preprocessing, the SVM used class weighting
with codebook-aware one-hot features, and the forest used an engineered
functional-limitation count. Differences may therefore reflect preprocessing
and imbalance strategies as well as the estimators themselves. A common-pipeline
ablation would provide a more controlled comparison.

\section{Limitations and Conclusion}

Thresholds were selected and reported on the same validation split, so the
local F1 values are tuning estimates rather than independent generalization
estimates. Repeated cross-validation or an untouched holdout would better
quantify uncertainty. BRFSS predictors and labels are self-reported, and their
substantial, potentially informative missingness can introduce bias. A linear
SVM may also miss nonlinear interactions, while a deeper forest increased
computational cost without clear gains.

Within these limitations, all three models benefited from threshold tuning.
The class-weighted linear SVM produced the highest local F1 and nearly identical
AIcrowd performance, making it the strongest current F1-oriented submission.
Future work should evaluate all models under identical preprocessing and use
repeated validation before fixing a final operating threshold.

\end{document}
